\documentclass[10pt,a4paper]{amsart}
\usepackage{amsmath,amssymb,amsthm,graphicx}
\usepackage{verbatim}
\newenvironment{reference}{\comment}{\endcomment}
%\newenvironment{reference}{}{}
\newenvironment{slogan}{\comment}{\endcomment}
\newenvironment{history}{\comment}{\endcomment}

% For commutative diagrams we use Xy-pic
\usepackage[all]{xy}

% We use 2cell for 2-commutative diagrams.
\xyoption{2cell}
% Theorem environments.
%
\theoremstyle{plain}
\newtheorem{theorem}[subsection]{Theorem}
\newtheorem{proposition}[subsection]{Proposition}
\newtheorem{lemma}[subsection]{Lemma}

\theoremstyle{definition}
\newtheorem{definition}[subsection]{Definition}
\newtheorem{example}[subsection]{Example}
\newtheorem{exercise}[subsection]{Exercise}
\newtheorem{situation}[subsection]{Situation}

\theoremstyle{remark}
\newtheorem{remark}[subsection]{Remark}
\newtheorem{remarks}[subsection]{Remarks}

\numberwithin{equation}{subsection}

% Macros
%
\def\lim{\mathop{\mathrm{lim}}\nolimits}
\def\colim{\mathop{\mathrm{colim}}\nolimits}
\def\Spec{\mathop{\mathrm{Spec}}}
\def\Hom{\mathop{\mathrm{Hom}}\nolimits}
\def\Ext{\mathop{\mathrm{Ext}}\nolimits}
\def\SheafHom{\mathop{\mathcal{H}\!\mathit{om}}\nolimits}
\def\SheafExt{\mathop{\mathcal{E}\!\mathit{xt}}\nolimits}
\def\Sch{\mathit{Sch}}
\def\Mor{\mathop{\mathrm{Mor}}\nolimits}
\def\Ob{\mathop{\mathrm{Ob}}\nolimits}
\def\Sh{\mathop{\mathit{Sh}}\nolimits}
\def\NL{\mathop{N\!L}\nolimits}
\def\CH{\mathop{\mathrm{CH}}\nolimits}
\def\proetale{{pro\text{-}\acute{e}tale}}
\def\etale{{\acute{e}tale}}
\def\QCoh{\mathit{QCoh}}
\def\Ker{\mathop{\mathrm{Ker}}}
\def\Im{\mathop{\mathrm{Im}}}
\def\Coker{\mathop{\mathrm{Coker}}}
\def\Coim{\mathop{\mathrm{Coim}}}

% Boxtimes
%
\DeclareMathSymbol{\boxtimes}{\mathbin}{AMSa}{"02}

%
% Macros for moduli stacks/spaces
%
\def\QCohstack{\mathcal{QC}\!\mathit{oh}}
\def\Cohstack{\mathcal{C}\!\mathit{oh}}
\def\Spacesstack{\mathcal{S}\!\mathit{paces}}
\def\Quotfunctor{\mathrm{Quot}}
\def\Hilbfunctor{\mathrm{Hilb}}
\def\Curvesstack{\mathcal{C}\!\mathit{urves}}
\def\Polarizedstack{\mathcal{P}\!\mathit{olarized}}
\def\Complexesstack{\mathcal{C}\!\mathit{omplexes}}
% \Pic is the operator that assigns to X its picard group, usage \Pic(X)
% \Picardstack_{X/B} denotes the Picard stack of X over B
% \Picardfunctor_{X/B} denotes the Picard functor of X over B
\def\Pic{\mathop{\mathrm{Pic}}\nolimits}
\def\Picardstack{\mathcal{P}\!\mathit{ic}}
\def\Picardfunctor{\mathrm{Pic}}
\def\Deformationcategory{\mathcal{D}\!\mathit{ef}}
\usepackage{cleveref}
\crefname{originallemma}{original lemma}{original lemmas}
\Crefname{originallemma}{Original lemma}{Original lemmas}
\makeatletter
\newlabel{lemma-lisse-etale}{{0788}{source}{}{sourcepointer-1}{}}
\newlabel{lemma-lisse-etale@cref}{{[originallemma][0][]0788}{[1][0][]source}{}{}{}}
\makeatother
\crefname{originalreference}{original reference}{original references}
\Crefname{originalreference}{Original reference}{Original references}
\makeatletter
\newlabel{sites-modules-definition-site-local}{{03DL}{source}{}{sourcepointer-2}{}}
\newlabel{sites-modules-definition-site-local@cref}{{[originalreference][0][]03DL}{[1][0][]source}{}{}{}}
\makeatother
\crefname{originalreference}{original reference}{original references}
\Crefname{originalreference}{Original reference}{Original references}
\makeatletter
\newlabel{sites-modules-definition-global}{{03DE}{source}{}{sourcepointer-3}{}}
\newlabel{sites-modules-definition-global@cref}{{[originalreference][0][]03DE}{[1][0][]source}{}{}{}}
\makeatother
\crefname{originalreference}{original reference}{original references}
\Crefname{originalreference}{Original reference}{Original references}
\makeatletter
\newlabel{stacks-sheaves-definition-pullback}{{06W1}{source}{}{sourcepointer-4}{}}
\newlabel{stacks-sheaves-definition-pullback@cref}{{[originalreference][0][]06W1}{[1][0][]source}{}{}{}}
\makeatother
\crefname{originallemma}{original lemma}{original lemmas}
\Crefname{originallemma}{Original lemma}{Original lemmas}
\makeatletter
\newlabel{lemma-lisse-etale-functorial}{{07AT}{source}{}{sourcepointer-5}{}}
\newlabel{lemma-lisse-etale-functorial@cref}{{[originallemma][0][]07AT}{[1][0][]source}{}{}{}}
\makeatother
\crefname{originallemma}{original lemma}{original lemmas}
\Crefname{originallemma}{Original lemma}{Original lemmas}
\makeatletter
\newlabel{lemma-check-qc-on-etale-covering}{{07AZ}{source}{}{sourcepointer-6}{}}
\newlabel{lemma-check-qc-on-etale-covering@cref}{{[originallemma][0][]07AZ}{[1][0][]source}{}{}{}}
\makeatother
\crefname{originallemma}{original lemma}{original lemmas}
\Crefname{originallemma}{Original lemma}{Original lemmas}
\makeatletter
\newlabel{lemma-lisse-etale-modules}{{0789}{source}{}{sourcepointer-7}{}}
\newlabel{lemma-lisse-etale-modules@cref}{{[originallemma][0][]0789}{[1][0][]source}{}{}{}}
\makeatother
\crefname{originalreference}{original reference}{original references}
\Crefname{originalreference}{Original reference}{Original references}
\makeatletter
\newlabel{categories-lemma-adjoint-exact}{{0038}{source}{}{sourcepointer-8}{}}
\newlabel{categories-lemma-adjoint-exact@cref}{{[originalreference][0][]0038}{[1][0][]source}{}{}{}}
\makeatother
\crefname{originallemma}{original lemma}{original lemmas}
\Crefname{originallemma}{Original lemma}{Original lemmas}
\makeatletter
\newlabel{lemma-lisse-etale-structure-sheaf}{{078A}{source}{}{sourcepointer-9}{}}
\newlabel{lemma-lisse-etale-structure-sheaf@cref}{{[originallemma][0][]078A}{[1][0][]source}{}{}{}}
\makeatother
\crefname{originalreference}{original reference}{original references}
\Crefname{originalreference}{Original reference}{Original references}
\makeatletter
\newlabel{stacks-sheaves-lemma-characterize-quasi-coherent-bis}{{0EM8}{source}{}{sourcepointer-10}{}}
\newlabel{stacks-sheaves-lemma-characterize-quasi-coherent-bis@cref}{{[originalreference][0][]0EM8}{[1][0][]source}{}{}{}}
\makeatother
\title{g\_! of a quasi-coherent module on the lisse-etale (flat-fppf) site is quasi-coherent}
\author{The Stacks Project Authors}
\date{stacks-project:stacks-cohomology.tex:local-snapshot-20261012; GNU FDL 1.2 or later}
\begin{document}
\maketitle
\noindent\textit{Editorial note.} All delivered mathematical text is the author's own, copied verbatim from the pinned Stacks Project chapter file stacks-cohomology.tex; no statement or proof step was rewritten, reordered or removed.\par\medskip
\noindent\textit{Editorial note.} The counted span cites the lemmas recorded in assembly.external\_references with canonical tags as navigation pointers.\par\medskip
\noindent\textit{Allowed standard prerequisites (editorial).}\begin{itemize}
\item Functoriality of the lisse-etale morphism of topoi. Cohomology of Algebraic Stacks, Lemma stacks-cohomology-lemma-lisse-etale-functorial. g\_\{i,!\}(f'\_i)\textasciicircum{}\{-1\} = f\_i\textasciicircum{}\{-1\}g\_!.
\item Descent of quasi-coherence along coverings. Cohomology of Algebraic Stacks, Lemma stacks-cohomology-lemma-check-qc-on-etale-covering. Reduces to the globally presented case.
\item Structure sheaf of the lisse-etale site. Cohomology of Algebraic Stacks, Lemma stacks-cohomology-lemma-lisse-etale-structure-sheaf. g\_!O = O\_X.
\end{itemize}\medskip
\ifdefined\hypertarget\hypertarget{sourcepointer-1}{}\fi
\noindent\textit{Source reference.} Original lemma 0788: sources/stacks-chow/tags.txt, source lines 9228--9228. Cross-chapter standard prerequisite cited by the author inside the counted span; it is used as a pointer and its own proof is not part of this item.\par\medskip
\ifdefined\hypertarget\hypertarget{sourcepointer-2}{}\fi
\noindent\textit{Source reference.} Original reference 03DL: sources/stacks-chow/tags.txt, source lines 4424--4424. Cross-chapter standard prerequisite cited by the author inside the counted span; it is used as a pointer and its own proof is not part of this item.\par\medskip
\ifdefined\hypertarget\hypertarget{sourcepointer-3}{}\fi
\noindent\textit{Source reference.} Original reference 03DE: sources/stacks-chow/tags.txt, source lines 4417--4417. Cross-chapter standard prerequisite cited by the author inside the counted span; it is used as a pointer and its own proof is not part of this item.\par\medskip
\ifdefined\hypertarget\hypertarget{sourcepointer-4}{}\fi
\noindent\textit{Source reference.} Original reference 06W1: sources/stacks-chow/tags.txt, source lines 8800--8800. Cross-chapter standard prerequisite cited by the author inside the counted span; it is used as a pointer and its own proof is not part of this item.\par\medskip
\ifdefined\hypertarget\hypertarget{sourcepointer-5}{}\fi
\noindent\textit{Source reference.} Original lemma 07AT: sources/stacks-chow/tags.txt, source lines 9318--9318. Cross-chapter standard prerequisite cited by the author inside the counted span; it is used as a pointer and its own proof is not part of this item.\par\medskip
\ifdefined\hypertarget\hypertarget{sourcepointer-6}{}\fi
\noindent\textit{Source reference.} Original lemma 07AZ: sources/stacks-chow/tags.txt, source lines 9324--9324. Cross-chapter standard prerequisite cited by the author inside the counted span; it is used as a pointer and its own proof is not part of this item.\par\medskip
\ifdefined\hypertarget\hypertarget{sourcepointer-7}{}\fi
\noindent\textit{Source reference.} Original lemma 0789: sources/stacks-chow/tags.txt, source lines 9229--9229. Cross-chapter standard prerequisite cited by the author inside the counted span; it is used as a pointer and its own proof is not part of this item.\par\medskip
\ifdefined\hypertarget\hypertarget{sourcepointer-8}{}\fi
\noindent\textit{Source reference.} Original reference 0038: sources/stacks-chow/tags.txt, source lines 125--125. Cross-chapter standard prerequisite cited by the author inside the counted span; it is used as a pointer and its own proof is not part of this item.\par\medskip
\ifdefined\hypertarget\hypertarget{sourcepointer-9}{}\fi
\noindent\textit{Source reference.} Original lemma 078A: sources/stacks-chow/tags.txt, source lines 9230--9230. Cross-chapter standard prerequisite cited by the author inside the counted span; it is used as a pointer and its own proof is not part of this item.\par\medskip
\ifdefined\hypertarget\hypertarget{sourcepointer-10}{}\fi
\noindent\textit{Source reference.} Original reference 0EM8: sources/stacks-chow/tags.txt, source lines 18298--18298. Cross-chapter standard prerequisite cited by the author inside the counted span; it is used as a pointer and its own proof is not part of this item.\par\medskip
\begin{lemma}
\label{lemma-shriek-quasi-coherent}
Let $\mathcal{X}$ be an algebraic stack. Notation as in
Lemma \ref{lemma-lisse-etale}.
\begin{enumerate}
\item Let $\mathcal{H}$ be a quasi-coherent
$\mathcal{O}_{\mathcal{X}_{lisse,\etale}}$-module 
on the lisse-\'etale site of $\mathcal{X}$. Then $g_!\mathcal{H}$ is a
quasi-coherent module on $\mathcal{X}$.
\item Let $\mathcal{H}$ be a quasi-coherent
$\mathcal{O}_{\mathcal{X}_{flat,fppf}}$-module 
on the flat-fppf site of $\mathcal{X}$. Then $g_!\mathcal{H}$ is a
quasi-coherent module on $\mathcal{X}$.
\end{enumerate}
\end{lemma}

\begin{proof}
Pick a scheme $U$ and a surjective smooth morphism $x : U \to \mathcal{X}$.
By
Modules on Sites, Definition \ref{sites-modules-definition-site-local}
there exists an \'etale (resp.\ fppf) covering
$\{U_i \to U\}_{i \in I}$ such that each pullback $f_i^{-1}\mathcal{H}$
has a global presentation (see
Modules on Sites, Definition \ref{sites-modules-definition-global}).
Here $f_i : U_i \to \mathcal{X}$ is the composition
$U_i \to U \to \mathcal{X}$ which is a morphism of algebraic stacks.
(Recall that the pullback ``is'' the restriction to $\mathcal{X}/f_i$, see
Sheaves on Stacks, Definition \ref{stacks-sheaves-definition-pullback}
and the discussion following.) Since each $f_i$ is smooth (resp.\ flat) by
Lemma \ref{lemma-lisse-etale-functorial}
we see that $f_i^{-1}g_!\mathcal{H} = g_{i, !}(f'_i)^{-1}\mathcal{H}$.
Using Lemma \ref{lemma-check-qc-on-etale-covering}
we reduce the statement of the lemma to the case where $\mathcal{H}$
has a global presentation. Say we have
$$
\bigoplus\nolimits_{j \in J} \mathcal{O} \longrightarrow
\bigoplus\nolimits_{i \in I} \mathcal{O} \longrightarrow
\mathcal{H} \longrightarrow 0
$$
of $\mathcal{O}$-modules where
$\mathcal{O} = \mathcal{O}_{\mathcal{X}_{lisse,\etale}}$
(resp.\ $\mathcal{O} = \mathcal{O}_{\mathcal{X}_{flat,fppf}}$).
Since $g_!$ commutes with arbitrary colimits (as a left adjoint functor, see
Lemma \ref{lemma-lisse-etale-modules} and
Categories, Lemma \ref{categories-lemma-adjoint-exact})
we conclude that there exists an exact sequence
$$
\bigoplus\nolimits_{j \in J} g_!\mathcal{O} \longrightarrow
\bigoplus\nolimits_{i \in I} g_!\mathcal{O} \longrightarrow
g_!\mathcal{H} \longrightarrow 0
$$
Lemma \ref{lemma-lisse-etale-structure-sheaf}
shows that $g_!\mathcal{O} = \mathcal{O}_\mathcal{X}$.
In case (2) we are done. In case (1) we apply
Sheaves on Stacks, Lemma
\ref{stacks-sheaves-lemma-characterize-quasi-coherent-bis}
to conclude.
\end{proof}

\end{document}
